\chapter*{Preface}
\addcontentsline{toc}{chapter}{Preface}

These notes grew out of the Models of Neural Systems course at the Bernstein Center for Computational Neuroscience in Berlin. The course bridges neurobiology and mathematics, building from the simplest models of neural computation up through the biophysics of ion channels, the dynamical systems theory of excitability, and the spatial and network models that connect single neurons to circuits.

The goal of these notes is not exhaustive coverage but \emph{deep understanding}. Each chapter is built around a small number of core ideas, developed from physical intuition through to mathematical formalism. Derivations are shown in full not because the algebra is the point, but because following the logic from assumptions to conclusions is how real understanding forms.

Throughout, we emphasize three recurring themes:
\begin{itemize}
  \item \textbf{What is the modeling goal?} Every equation exists to answer a question about the nervous system.
  \item \textbf{What are the key variables and parameters?} Identifying state variables and understanding what each parameter controls is the foundation of model intuition.
  \item \textbf{What qualitative behavior does the model explain?} Thresholds, tuning curves, oscillations, propagation, memory---the real payoff of mathematical modeling is explaining \emph{why} these phenomena occur.
\end{itemize}

These notes are meant to be read actively. Work through the derivations with pen and paper. Attempt the exercises before looking at solutions. When you encounter a key equation, ask yourself: what does each term do? What happens if I change a parameter? What breaks if I remove a term?

\bigskip
\noindent\textit{Berlin, February 2026}
